\chapter{Ear Barotrauma}

\section*{Definition}
Ear barotrauma (barotitis media) is tissue damage resulting from failure of the Eustachian tube to equalize pressure across the tympanic membrane when ambient pressure changes, most commonly during aircraft descent or scuba diving ascent.

\section*{Pathophysiology}
When the Eustachian tube cannot adequately open (due to edema, obstruction, or inadequate technique), the negative middle ear pressure during descent causes inward retraction of the tympanic membrane, mucosal congestion, serosanguineous effusion, and potentially tympanic membrane rupture if pressure exceeds 100--150 mmHg.

\section*{Etiology}
\begin{itemize}
  \item \textbf{Rapid pressure changes:} Airplane descent, scuba diving ascent/descent, hyperbaric chamber
  \item \textbf{Eustachian tube dysfunction:} Upper respiratory infection, allergic rhinitis, nasal congestion, sinusitis
  \item \textbf{Anatomical obstruction:} Nasal polyps, adenoid hypertrophy, cleft palate, nasopharyngeal tumor
  \item \textbf{Inadequate equalization technique:} Failure to yawn, swallow, or perform valsalva during pressure changes
\end{itemize}

\section*{Indications for Evaluation}
Seek care if:
\begin{itemize}
  \item Severe ear pain during or after pressure change that does not resolve within hours
  \item Hearing loss, ear fullness, or aural fullness persisting beyond 24 hours
  \item Tympanic membrane rupture (sudden pain relief followed by bloody otorrhea)
  \item Vertigo or tinnitus after diving (suspect inner ear barotrauma or perilymphatic fistula)
  \item Recurrent barotrauma with minimal pressure changes
\end{itemize}

\section*{Related Symptoms}
\begin{itemize}
  \item Otalgia (ear pain)
  \item Aural fullness and pressure sensation
  \item Conductive hearing loss
  \item Tinnitus
  \item Vertigo (if inner ear involved)
\end{itemize}
\textbf{Note:} Scuba divers with inner ear barotrauma require immediate ENT evaluation and may need surgical exploration for perilymph fistula; flying should be avoided until symptoms resolve.

\section*{Self-Care / Home Management}
\begin{itemize}
  \item Active Eustachian tube opening maneuvers: yawning, swallowing, jaw movement, or gentle valsalva
  \item Decongestants (oral pseudoephedrine or topical oxymetazoline) before anticipated pressure changes
  \item Oral antihistamines if allergic component suspected
  \item Earplugs designed for pressure equalization during flying
  \item Avoid flying or diving with active nasal congestion or upper respiratory infection
  \item Benzodiazepines or antiemetics for associated vertigo
\end{itemize}

\section*{Diagnostic Approach}
\begin{itemize}
  \item Otoscopy: Tympanic membrane retraction, hemotympanum, middle ear effusion, or perforation
  \item Pneumatic otoscopy to assess tympanic membrane mobility
  \item Pure-tone audiometry to quantify conductive hearing loss
  \item Tympanometry demonstrating flat curve (type B) in effusion
  \item Tuning fork tests (Rinne and Weber) to distinguish conductive from sensorineural loss
  \item CT temporal bone if suspected inner ear barotrauma or perilymph fistula
\end{itemize}

\section*{Treatment / Management Overview}
\begin{itemize}
  \item Mild barotrauma (TM retraction, no effusion): Observation, decongestants, and autoinflation
  \item Middle ear effusion: Oral decongestants and corticosteroids; symptoms typically resolve in days to weeks
  \item Hemotympanum: Conservative management with avoidance of pressure changes; resolution in 2--4 weeks
  \item Tympanic membrane perforation: Keep ear dry, avoid water exposure, topical antibiotics if infection; most heal spontaneously in 4--6 weeks
  \item Incision and drainage of middle ear effusion (myringotomy) if persistent beyond 3 months
  \item Pressure equalization tube (grommet) insertion for recurrent barotrauma or chronic ET dysfunction
\end{itemize}

\clearpage
